\documentclass[a4paper]{article}

\usepackage[utf8]{inputenc}
\usepackage[spanish]{babel}
\usepackage{amsmath,amssymb,amsfonts}
\usepackage{enumitem}
\usepackage{graphicx}
\usepackage{float}

% Definir los entornos necesarios para exams
\newenvironment{question}{}{}
\newenvironment{solution}{}{}
\newenvironment{answerlist}{\begin{enumerate}}{\end{enumerate}}

\usepackage{Sweave}
\begin{document}
\Sconcordance{concordance:consumo_gas_natural_porcentaje_maximo_aleatorio_interpretacion_representacion.tex:consumo_gas_natural_porcentaje_maximo_aleatorio_interpretacion_representacion.Rnw:1 %
14 1 1 0 2 1 1 207 1 84 4 1 1 10 2 1 1 2 4 0 1 2 17 1 1 7 5 0 1 3 23 1 %
1 17 1 0 1 3 14 1 1 2 4 0 1 2 17 1}




\begin{question}

Valentina y Francisco viven en un apartamento y comparten el pago de los gastos. En octubre consumieron 12 metros cúbicos de gas natural y la factura fue de \$14.000, incluido el cargo fijo de \$3.200 que es el mismo todos los meses. La gráfica muestra el consumo histórico en metros cúbicos de ese servicio.

\includegraphics[width=0.8\textwidth]{grafico_consumo_gas}
Un hogar puede consumir máximo 25 metros cúbicos en un mes. ¿A qué porcentaje del consumo máximo posible corresponde el consumo de mayo?

\begin{answerlist}
  \item 19%
  \item 83%
  \item 56%
  \item 76%
\end{answerlist}\end{question}

\begin{solution}

\subsection*{Análisis del problema}

Este problema requiere \textbf{interpretación de gráficos de barras} y \textbf{cálculo de porcentajes}. Los pasos son:

\begin{enumerate}
\item \textbf{Leer correctamente la gráfica} para identificar el consumo de mayo
\item \textbf{Identificar el consumo máximo posible} según el enunciado
\item \textbf{Aplicar la fórmula de porcentaje} para encontrar la relación
\end{enumerate}

\subsection*{Paso 1: Lectura de la gráfica}

Observando la gráfica de consumo histórico de Valentina y Francisco:

\textbf{Consumos mensuales registrados:}\begin{itemize}\item Abril: 14 metros cúbicos
\item Mayo: 19 metros cúbicos $\leftarrow$ \textbf{MES DE LA PREGUNTA}
\item Junio: 14 metros cúbicos
\item Julio: 14 metros cúbicos
\item Agosto: 15 metros cúbicos
\item Septiembre: 14 metros cúbicos
\item Octubre: 12 metros cúbicos $\leftarrow$ \textbf{MES DE LA FACTURA}\end{itemize}
\textbf{El consumo de mayo fue de 19 metros cúbicos.}

\subsection*{Paso 2: Datos del problema}

\begin{itemize}
\item \textbf{Consumo máximo posible:} 25 metros cúbicos (dato del enunciado)
\item \textbf{Consumo real de mayo:} 19 metros cúbicos (dato de la gráfica)
\end{itemize}

\subsection*{Paso 3: Aplicación de la fórmula de porcentaje}

Para calcular qué porcentaje representa el consumo de mayo respecto al máximo posible:

$$\text{Porcentaje} = \frac{\text{Valor observado}}{\text{Valor máximo}} \times 100\%$$

Sustituyendo los valores:

$$\text{Porcentaje} = \frac{19}{25} \times 100\% = 76\%$$

\subsection*{Verificación y análisis de distractores}

\textbf{Respuesta correcta:} 76\%

\textbf{Análisis de errores conceptuales comunes:}\begin{itemize}\item \textbf{19\%}: Error conceptual - confundir el valor absoluto (19 m³) con el porcentaje
\item \textbf{56\%}: Error de lectura - usar el consumo de abril en lugar de mayo
\item \textbf{83\%}: Error de cálculo - aplicar incorrectamente la fórmula de porcentaje\end{itemize}
\subsection*{Conclusión}

El consumo de mayo (19 metros cúbicos) representa el \textbf{76\%} del consumo máximo posible (25 metros cúbicos).

Esta respuesta es coherente porque:

\begin{itemize}
\item Se basa en una lectura correcta de la gráfica
\item Aplica correctamente la fórmula de porcentaje
\item El resultado está dentro del rango esperado (0\% a 100\%)
\end{itemize}

\textbf{Verificación adicional}: La factura de octubre (\$14.000) corresponde exactamente al consumo mostrado en la gráfica (12 m³ $\times$ \$900 + \$3.200 cargo fijo).

\begin{answerlist}
  \item Incorrecto
  \item Incorrecto
  \item Incorrecto
  \item Correcto
\end{answerlist}\end{solution}

%% META-INFORMATION
%% \exname{Consumo Gas Natural Porcentaje Máximo}
%% \extype{schoice}
%% \exsolution{0001}
%% \exshuffle{TRUE}
%% \exsection{Estadística|Interpretación de gráficos|Porcentajes|Análisis de datos}
%% \exextra[Type]{Cálculo}
%% \exextra[Program]{R}
%% \exextra[Language]{es}
%% \exextra[Level]{2}
%% \exextra[Competencia]{Interpretación y representación}
%% \exextra[Componente]{Aleatorio y sistemas de datos}
%% \exextra[Contexto]{Familiar}
%% \exextra[Dificultad]{Media}

\end{document}
